% ----------------------------------------------------------------------
%  3.7.3. Теплопроводность твёрдых тел
% ----------------------------------------------------------------------
\subsection{Теплопроводность}\label{ssec:thermal-cond}

\paragraph{Основные определения}
Теплопроводность~— перенос энергии теплового движения от горячих
участков тела к холодным. Она описывается тем же типом уравнения, что и
перенос заряда (закон Ома) и вещества (первый закон Фика,
п.~\ref{ssec:fick}): плотность потока тепла $q$~— количество теплоты,
проходящее за 1~с через 1~м$^2$,~— пропорциональна градиенту температуры
(закон Фурье):
\begin{keyeq}\label{eq:fourier}
q=-\lambda\,\frac{\d T}{\d x},\qquad a_T=\frac{\lambda}{\rho c} .
\end{keyeq}
Коэффициент $\lambda$, Вт/(м$\cdot$К),~— \term{теплопроводность}; знак
«минус» означает, что тепло течёт в сторону понижения температуры. Для
нестационарных процессов важна \term{температуропроводность} $a_T$, м$^2$/с
($\rho$~— плотность, $c$~— удельная теплоёмкость): она входит в уравнение
$\partial T/\partial t=a_T\,\partial^2T/\partial x^2$ (аналог второго закона
Фика) и определяет, как быстро выравнивается температура.

\paragraph{Кто переносит тепло}
\begin{itemize}
\item В \emph{диэлектриках} (ионных и ковалентных кристаллах)~— только
фононы; теплопроводность обычно на 1--2 порядка ниже, чем у металлов
(с важными исключениями~— алмаз, BeO, SiC).
\item В \emph{металлах}~— электроны проводимости ($\lambda_\text{эл}$) и
фононы ($\lambda_\text{реш}$). В чистых металлах главный механизм~—
электронный, решёточная доля~— лишь несколько процентов; в сплавах обе
составляющие могут быть сравнимы.
\item В \emph{полупроводниках} к решёточной теплопроводности добавляется
перенос тепла носителями заряда (и экситонами).
\end{itemize}

\subsubsection{Теплопроводность решётки}

\paragraph{Без ангармонизма теплопроводность была бы бесконечной}
Если бы атомы колебались строго гармонически, нормальные колебания были бы
независимыми и распространялись бы по кристаллу, не рассеиваясь. Тогда
тепловой импульс, созданный на горячем конце, бежал бы по кристаллу со
скоростью звука, ничего не теряя, и даже при бесконечно малой разности
температур поток тепла мог бы быть сколь угодно большим.

Реальные колебания ангармоничны (член $bx^3$ в разложении потенциальной
энергии, п.~\ref{ssec:heat-capacity}). Поэтому фононы
\emph{взаимодействуют}: фонон может распасться на два, два фонона~—
слиться в один. Эти процессы \term{фонон-фононного рассеяния}
перераспределяют энергию между колебаниями разных частот и
устанавливают тепловое равновесие. Вероятность рассеяния характеризуют
сечением $\sigma_\text{ф}$, и так как рассеяние вызывается ангармонизмом,
естественно считать $\sigma_\text{ф}\propto b^2$.

\paragraph{Фононный газ}
По кинетической теории газов теплопроводность $\lambda=\frac13 c_V\bar\upsilon l$,
где $c_V$~— теплоёмкость единицы объёма, $\bar\upsilon$~— скорость частиц,
$l$~— длина свободного пробега. Для фононного газа скорость~— это скорость
звука, а длина пробега $l_\text{ф}\approx1/(n_\text{ф}\sigma_\text{ф})$
обратно пропорциональна концентрации фононов $n_\text{ф}$:
\begin{keyeq}\label{eq:lambda-lattice}
\lambda_\text{реш}=\frac13\,c_V\,\upsilon_\text{зв}\,l_\text{ф},
\qquad l_\text{ф}\propto\frac{1}{n_\text{ф}b^2},\qquad
n_\text{ф}\propto\frac{1}{e^{\theta_D/T}-1} .
\end{keyeq}

\begin{table}[ht]
\centering
\caption{Тепловые свойства кристаллов кварца и NaCl}\label{tab:phonon-l}
\small
\begin{tabular}{lcccc}
\toprule
Кристалл & $T$, К & $c_V$, $10^6$~Дж/(м$^3\cdot$К) & $\lambda$, Вт/(м$\cdot$К) & $l_\text{ф}$, нм\\
\midrule
SiO$_2$ & 270 & 2,0 & 13 & $\approx4$\\
        &  80 & 0,5 & 50 & $\approx60$\\
NaCl    & 270 & 1,9 &  7 & $\approx4$\\
        &  80 & 1,0 & 26 & $\approx25$\\
\bottomrule
\multicolumn{5}{l}{\footnotesize $l_\text{ф}=3\lambda/(c_V\upsilon_\text{зв})$ при $\upsilon_\text{зв}\approx5$~км/с (кварц) и $3$~км/с (NaCl).}
\end{tabular}
\end{table}

Порядок величин даёт табл.~\ref{tab:phonon-l}: при комнатной температуре
фонон пролетает между столкновениями лишь несколько нанометров (порядка
десяти межатомных расстояний), а при охлаждении длина пробега быстро
растёт\ispr{В оригинале (табл.~3.25) длины пробега фононов завышены
примерно в 100 раз (400, 5400, 210, 1500~нм); значения в таблице
пересчитаны по формуле $l_\text{ф}=3\lambda/(c_V\upsilon_\text{зв})$ из
приведённых там же $c_V$ и $\lambda$.}.

\begin{masshtab}[фононы]{-9.3/-8.9/{длина волны тепловых фононов}, -9/-8.3/{пробег при 300\,К}, -4/-2/{пробег при $T\ll\theta_D$}}
Тепловые фононы при комнатной температуре~— волны длиной около
нанометра (длина волны $\sim h\upsilon_\text{зв}/kT$), а пролетают они между
столкновениями лишь несколько нанометров. При глубоком охлаждении
фононов становится мало, пробег растёт и в совершенном кристалле
доходит до размеров образца; тогда теплопроводность зависит от размеров
и состояния поверхности образца. Ультразвук и звук в том же кристалле~—
тоже фононы, но с длинами волн от микрометров до метров.
\end{masshtab}

\paragraph{Высокие температуры ($T>\theta_D$)}
Здесь $e^{\theta_D/T}-1\approx\theta_D/T$, так что $n_\text{ф}\propto T$, а
$c_V\approx\text{const}$. Следовательно,
\begin{equation}\label{eq:lambda-highT}
\lambda_\text{реш}\propto\frac{c_V\upsilon_\text{зв}}{T\,b^2}\propto\frac1T .
\end{equation}
Формула объясняет и связь теплопроводности с прочностью связи: при
ослаблении связи скорость звука падает, а ангармонизм растёт (амплитуды
колебаний при той же температуре больше), и оба фактора уменьшают
$\lambda_\text{реш}$. Это хорошо видно на элементах со структурой алмаза
(табл.~\ref{tab:diamond-lambda}). Кроме того, $\lambda_\text{реш}$ быстро
уменьшается с ростом массы атомов: у лёгких элементов верхней части
таблицы Менделеева (B, C, Si) она составляет десятки и сотни
Вт/(м$\cdot$К), в середине таблицы~— единицы, у тяжёлых элементов~—
десятые доли Вт/(м$\cdot$К).

\begin{table}[ht]
\centering
\caption{Теплота сублимации и теплопроводность элементов со структурой
алмаза (300~К)}\label{tab:diamond-lambda}
\small
\begin{tabular}{lcc}
\toprule
Вещество & $Q_\text{субл}$, кДж/моль & $\lambda_\text{реш}$, Вт/(м$\cdot$К)\\
\midrule
Алмаз & 712 & $\approx2000$\\
Кремний & 461 & 150\\
Германий & 370 & 60\\
\bottomrule
\end{tabular}
\end{table}

У кристаллов из лёгких атомов с жёсткими связями решёточная
теплопроводность бывает огромной: у алмаза при комнатной температуре она
примерно впятеро выше, чем у самого теплопроводного металла~— серебра
($\lambda_\text{Ag}\approx430$~Вт/(м$\cdot$К))\ispr{В оригинале
(табл.~3.26) теплоты сублимации приведены в $10^5$~Дж/моль (71,2; 46,1;
37,0), что в 10 раз больше правильных значений, а теплопроводность алмаза
указана 550~Вт/(м$\cdot$К); для чистых кристаллов алмаза она составляет
2000--2200~Вт/(м$\cdot$К), для кремния~— около 150, германия~— около
60~Вт/(м$\cdot$К).}.

\paragraph{Низкие температуры ($T\ll\theta_D$)}
Концентрация фононов резко падает, и длина пробега растёт. При
$T\lesssim\theta_D/20$ она становится сравнимой с размером кристалла~$L$ и
дальше расти не может: фононы рассеиваются на поверхности образца. Тогда
$l_\text{ф}\approx L=\text{const}$, и теплопроводность повторяет
температурную зависимость теплоёмкости:
\begin{equation}\label{eq:lambda-lowT}
\lambda_\text{реш}\propto c_V\propto T^3 .
\end{equation}

\paragraph{Промежуточная область}
При $T$ несколько ниже $\theta_D$ длина пробега растёт с охлаждением
\emph{экспоненциально}: $l_\text{ф}\propto e^{\theta_D/bT}$, где $b\approx2$--3,
и $\lambda_\text{реш}\propto c_V(T)\,e^{\theta_D/bT}$\ispr{В оригинале знак
показателя экспоненты обратный: $l_\text{ф}\sim e^{-\theta_D/T}$, т.\,е. длина
пробега убывала бы при охлаждении, что противоречит сказанному выше.}.
Причина в том, что тепловое сопротивление создают не все столкновения
фононов, а только те, в которых суммарный квазиимпульс выходит за пределы
зоны Бриллюэна (см. «Подробнее» ниже). Для этого нужны фононы с достаточно
большой энергией, а их число убывает как $e^{-\theta_D/bT}$.

\begin{glavnoe}
Теплопроводность диэлектрика проходит через \vazhno{максимум}
(рис.~\ref{fig:kappa-T}\,\textit{а}). При нагреве от абсолютного нуля растёт
число фононов, переносящих тепло ($\lambda\propto T^3$), но одновременно
усиливается их рассеяние друг на друге, и длина пробега падает. Сначала
побеждает первый фактор, потом второй, и при высоких температурах
$\lambda\propto1/T$. Положение и высота максимума зависят от размеров и
совершенства кристалла.
\end{glavnoe}

\begin{podrobnee}[Нормальные процессы и процессы переброса]
При столкновении фононов сохраняется энергия, а квазиимпульс~— с
точностью до вектора обратной решётки:
$\vb k_1+\vb k_2=\vb k_3+\vb G$. Если $\vb G=0$ (\emph{нормальный процесс}),
суммарный импульс фононного газа не меняется, и поток тепла не
затухает~— такие столкновения сопротивления не создают. Тепловое
сопротивление дают лишь \emph{процессы переброса} (Umklapp, $\vb G\ne0$), в
которых сумма $\vb k_1+\vb k_2$ выходит за зону Бриллюэна; для этого каждый
из сталкивающихся фононов должен иметь волновой вектор порядка половины
размера зоны и энергию порядка $k\theta_D/2$. Число таких фононов и
определяет множитель $e^{-\theta_D/bT}$ (Р.~Пайерлс, 1929).
\end{podrobnee}

В \emph{аморфных} диэлектриках (стёклах) область правильной структуры имеет
атомные размеры, рассеяние на нарушениях структуры преобладает при всех
температурах, и длина пробега~— порядка межатомного расстояния. Поэтому
теплопроводность стёкол мала (около 1~Вт/(м$\cdot$К) у кварцевого стекла
против 10--13 у кристаллического кварца) и при высоких температурах почти
не зависит от $T$, следуя за теплоёмкостью\ispr{В оригинале сказано, что при
низких температурах теплопроводность стёкол пропорциональна $T^3$, «что и
наблюдается на опыте». В действительности ниже ${\sim}1$~К у стёкол
$\lambda\propto T^2$, а при 5--20~К наблюдается характерное «плато»
(Р.~Зеллер, Р.~Пол, 1971); простая модель с постоянной длиной пробега
описывает лишь малую величину $\lambda$ и её слабую зависимость от $T$ при
высоких температурах.}.

Заметим, что теория хорошо описывает основные закономерности решёточной
теплопроводности, но точного значения $\lambda_\text{реш}$ предсказать пока не
позволяет.

\subsubsection{Теплопроводность металлов}

В металле тепло переносят и фононы, и электроны:
$\lambda=\lambda_\text{реш}+\lambda_\text{эл}$. Для электронного газа в формулу
кинетической теории подставим теплоёмкость электронов~\eqref{eq:C-el}
(на единицу объёма $c_\text{эл}=\frac{\pi^2}{2}nk\,\frac{kT}{\EF}$), скорость
Ферми $\upsilon_\text{F}$ и длину пробега электронов $l_\text{эл}$:
\begin{keyeq}\label{eq:lambda-el}
\lambda_\text{эл}=\frac13\,c_\text{эл}\upsilon_\text{F}l_\text{эл}
=\frac{\pi^2nk^2T}{3m\upsilon_\text{F}}\,l_\text{эл}
\end{keyeq}
($n$~— концентрация электронов; использовано $\EF=m\upsilon_\text{F}^2/2$)\ispr{В
оригинале в числителе стоит $\pi^3$; правильно $\pi^2$, как следует из
$c_\text{эл}=\pi^2nk^2T/(m\upsilon_\text{F}^2)$.}.

\paragraph{Температурная зависимость}
В формуле~\eqref{eq:lambda-el} от температуры зависят множитель $T$ и длина
пробега $l_\text{эл}$, которая определяется рассеянием электронов на фононах
и на дефектах:
\begin{itemize}
\item $T\gtrsim\theta_D$: $l_\text{эл}\propto1/n_\text{ф}\propto1/T$, и
$\lambda_\text{эл}=\text{const}$. У меди теплопроводность практически
постоянна выше 80--100~К (рис.~\ref{fig:kappa-T}\,\textit{б});
\item $T\ll\theta_D$: $n_\text{ф}\propto T^3$, $l_\text{эл}\propto T^{-3}$, и
$\lambda_\text{эл}\propto T^{-2}$;
\item вблизи 0~К фононов почти нет, электроны рассеиваются на примесях,
которые есть в любом металле: $l_\text{эл}\propto1/N_\text{пр}$ не зависит от
$T$, и $\lambda_\text{эл}\propto T$.
\end{itemize}
Тепловые сопротивления от разных механизмов рассеяния складываются (как
электросопротивления в правиле Маттиссена), поэтому при $T\ll\theta_D$
\begin{equation}\label{eq:W-metal}
\frac1{\lambda_\text{эл}}=aT^2+\frac bT ,
\end{equation}
где первый член~— рассеяние на фононах, второй~— на примесях и дефектах.
Отсюда максимум теплопроводности при $T_{\max}=(b/2a)^{1/3}$: чем грязнее
металл (больше $b$), тем выше температура максимума и тем ниже его высота.

\begin{figure}[ht]
\centering
\begin{subfigure}[b]{0.49\textwidth}\centering
\begin{tikzpicture}
\begin{axis}[width=0.98\linewidth, height=5.6cm, xmode=log, ymode=log,
  xmin=1, xmax=600, ymin=1, ymax=2e4,
  xlabel={$T$, К}, ylabel={$\lambda$, Вт/(м$\cdot$К)},
  legend style={at={(0.98,0.98)}, anchor=north east, font=\scriptsize}]
  \addplot[fmBlue, thick] table[x=T, y=big] {figures/th-kappa-diel.dat};
  \addlegendentry{$L=10$~мм}
  \addplot[fmRed, thick, dashed] table[x=T, y=small] {figures/th-kappa-diel.dat};
  \addlegendentry{$L=1$~мм}
  \node[font=\scriptsize, rotate=52] at (axis cs:2.6,700) {$\propto T^3$};
  \node[font=\scriptsize, anchor=south west] at (axis cs:150,90) {$\propto 1/T$};
\end{axis}
\end{tikzpicture}
\caption{}
\end{subfigure}\hfill
\begin{subfigure}[b]{0.49\textwidth}\centering
\begin{tikzpicture}
\begin{axis}[width=0.98\linewidth, height=5.6cm, xmode=log, ymode=log,
  xmin=1, xmax=400, ymin=150, ymax=1e4,
  xlabel={$T$, К}, ylabel={$\lambda$, Вт/(м$\cdot$К)},
  legend style={at={(0.98,0.03)}, anchor=south east, font=\scriptsize}]
  \addplot[fmBlue, thick] table[x=T, y=pure] {figures/th-kappa-cu.dat};
  \addlegendentry{чистый металл}
  \addplot[fmRed, thick, dashed] table[x=T, y=dirty] {figures/th-kappa-cu.dat};
  \addlegendentry{в 10 раз больше примесей}
  \node[font=\scriptsize, rotate=28] at (axis cs:1.6,3600) {$\propto T$};
  \node[font=\scriptsize, rotate=-50] at (axis cs:30,3000) {$\propto T^{-2}$};
  \node[font=\scriptsize, anchor=south] at (axis cs:180,480) {$\approx\text{const}$};
\end{axis}
\end{tikzpicture}
\caption{}
\end{subfigure}
\caption{Температурная зависимость теплопроводности (модельный расчёт):
\textit{а}~— диэлектрик с параметрами сапфира (кристаллы разного размера);
$\lambda=\frac13c_V\upsilon l$, где длину пробега ограничивают размер кристалла,
точечные дефекты и процессы переброса; \textit{б}~— медь двух степеней чистоты
по формуле~\eqref{eq:W-metal}, дополненной насыщением фононного
сопротивления при $T\gtrsim\theta_D/4$}
\label{fig:kappa-T}
\end{figure}

\paragraph{Порядок величины}
Для металлов $c_\text{эл}\approx0{,}01\,c_V\approx3\cdot10^4$~Дж/(м$^3\cdot$К),
$\upsilon_\text{F}\approx10^6$~м/с, $l_\text{эл}\approx10^{-8}$~м при комнатной
температуре. Формула~\eqref{eq:lambda-el} даёт
$\lambda_\text{эл}\approx\frac13\cdot3\cdot10^4\cdot10^6\cdot10^{-8}=10^2$~Вт/(м$\cdot$К)\ispr{В
оригинале в этой оценке стоит скорость звука; в формулу входит скорость
электронов на поверхности Ферми $\upsilon_\text{F}\approx10^6$~м/с.}~—
сотни ватт на метр-кельвин, в согласии с опытом (табл.~\ref{tab:lambda-metals})\ispr{Значения
табл.~\ref{tab:lambda-metals} даны по современным справочникам; в оригинале
(табл.~3.27): серебро~403, медь~384, золото~296, алюминий~210,
никель~60~Вт/(м$\cdot$К)~— для никеля это заметно заниженная величина.}.

\begin{table}[ht]
\centering
\caption{Теплопроводность металлов и сплава при 300~К}\label{tab:lambda-metals}
\small
\begin{tabular}{lc@{\qquad}lc}
\toprule
Металл & $\lambda$, Вт/(м$\cdot$К) & Металл (сплав) & $\lambda$, Вт/(м$\cdot$К)\\
\midrule
Серебро & 429 & Алюминий & 237\\
Медь & 401 & Никель & 91\\
Золото & 318 & Константан (Cu--40\,\%\,Ni) & 22\\
\bottomrule
\end{tabular}
\end{table}

\paragraph{Доля решётки в металле}
Отношение~\eqref{eq:lambda-lattice} к~\eqref{eq:lambda-el}:
\[
\frac{\lambda_\text{реш}}{\lambda_\text{эл}}=\frac{c_V}{c_\text{эл}}\cdot
\frac{\upsilon_\text{зв}}{\upsilon_\text{F}}\cdot\frac{l_\text{ф}}{l_\text{эл}}
\approx100\cdot\frac{5\cdot10^3}{10^6}\cdot\frac{10^{-9}}{10^{-8}}=5\cdot10^{-2}.
\]
У чистых металлов на решётку приходится лишь несколько процентов.
Но в \emph{сплавах} электроны рассеиваются в основном на атомах примеси,
$l_\text{эл}\propto1/N_\text{пр}$, и при большой концентрации длина пробега
электронов может сравняться с фононной: тогда
$\lambda_\text{эл}\approx\lambda_\text{реш}$. Так, у константана теплопроводность
во много раз ниже, чем у меди и никеля: электроны в нём рассеиваются
главным образом на искажениях решётки вокруг атомов другого сорта. А у
карбида TiC, имеющего металлическую проводимость, решёточная составляющая
даже преобладает над электронной ниже 310~К.

Теплопроводность металлов, как и электропроводность, меняется в магнитном
поле; наибольшие изменения (в 2--5 раз) наблюдаются в поперечном поле в
области максимума теплопроводности.

\paragraph{Закон Видемана--Франца--Лоренца}
В ряду металлов теплопроводность тем больше, чем больше
электропроводность. Г.~Видеман и Р.~Франц установили это на опыте, а
Л.~Лоренц показал, что отношение $\lambda/\sigma$ пропорционально абсолютной
температуре. Это сразу следует из электронной теории: деля~\eqref{eq:lambda-el}
на $\sigma=ne^2l_\text{эл}/(m\upsilon_\text{F})$~\eqref{eq:sigma-quantum}, получаем
(длина пробега и масса сокращаются)
\begin{keyeq}\label{eq:WF}
\frac{\lambda}{\sigma}=LT,\qquad L=\frac{\pi^2}{3}\left(\frac ke\right)^2=2{,}45\cdot10^{-8}~\text{Вт}\cdot\text{Ом}/\text{К}^2 .
\end{keyeq}
Постоянная $L$~— \term{число Лоренца}. Для большинства технически чистых
металлов при комнатной температуре опытные $L=(2{,}1$--$2{,}8)\cdot10^{-8}$~Вт$\cdot$Ом/К$^2$;
у ферромагнитных металлов $L$ аномально велико (у железа
$3{,}0\cdot10^{-8}$). В полупроводниках с невырожденным электронным газом
электронная составляющая тоже подчиняется закону, но с $L\approx2(k/e)^2$, а
значительную долю теплопроводности даёт решётка.

\begin{primer}[Теплопроводность меди по её электропроводности]
Удельная проводимость меди при 293~К $\sigma=5{,}96\cdot10^7$~См/м. По
закону~\eqref{eq:WF} $\lambda\approx2{,}45\cdot10^{-8}\cdot293\cdot5{,}96\cdot10^7\approx430$~Вт/(м$\cdot$К);
опыт~— 401~Вт/(м$\cdot$К), т.\,е. для меди $L\approx2{,}3\cdot10^{-8}$~Вт$\cdot$Ом/К$^2$.
\end{primer}

\begin{vnimanie}[Когда закон нарушается]
Закон точен, если электроны рассеиваются \emph{упруго}: при высоких
температурах ($T\gtrsim\theta_D$) и при самых низких, где главное~—
рассеяние на примесях. В промежуточной области рассеяние на фононах
неупругое, тепловой поток затухает быстрее электрического тока, и $L$
оказывается меньше теоретического. При высоких температурах в сплавах
отклонения даёт заметная решёточная теплопроводность.
\end{vnimanie}

Несмотря на приближённость, закон Видемана--Франца--Лоренца позволяет
переносить закономерности электропроводности на теплопроводность и
оценивать $\lambda$ по $\sigma$. Это важно на практике: тепло\-проводность
измерять гораздо труднее (мешают теплопотери), чем электропроводность.

\paragraph{Теплопроводность сплавов}
По аналогии с электропроводностью (п.~\ref{ssec:rho-alloys}):
\begin{itemize}
\item в непрерывном ряду неупорядоченных твёрдых растворов теплопроводность
падает по мере удаления от чистых компонентов; минимум, как правило, около
50~ат.\,\% (рис.~\ref{fig:lambda-alloys}\,\textit{а}). Сильнее всего
теплопроводность чистого металла снижают первые доли процента примеси,
дальнейшее легирование влияет слабее;
\item в гетерогенной (двухфазной) области теплопроводность меняется
примерно линейно с объёмной долей фаз между значениями
$\lambda_\alpha$ и $\lambda_\beta$ для фаз предельной концентрации
(рис.~\ref{fig:lambda-alloys}\,\textit{б}); для любого сплава из этой
области она лежит между значениями для фаз;
\item в композитах (волокна или пластины одного металла в матрице другого)
важна геометрия: при тепловом потоке вдоль слоёв
$\lambda_\parallel=\sum v_i\lambda_i$, поперёк~— $1/\lambda_\perp=\sum v_i/\lambda_i$
($v_i$~— объёмные доли).
\end{itemize}

\begin{figure}[ht]
\centering
\begin{subfigure}[b]{0.49\textwidth}\centering
\begin{tikzpicture}
\begin{axis}[width=0.92\linewidth, height=5.2cm, xmin=0, xmax=100, ymin=0, ymax=500,
  xlabel={Au, ат.\,\%}, ylabel={$\lambda$, Вт/(м$\cdot$К)}]
  \addplot[fmBlue, thick, domain=0:100, samples=101]
    {735/(1.59*(1-x/100)+2.21*x/100+34*(x/100)*(1-x/100))};
  \node[font=\scriptsize, anchor=west] at (axis cs:2,460) {Ag};
  \node[font=\scriptsize, anchor=east] at (axis cs:98,345) {Au};
\end{axis}
\end{tikzpicture}
\caption{}
\end{subfigure}\hfill
\begin{subfigure}[b]{0.49\textwidth}\centering
\begin{tikzpicture}
\begin{axis}[width=0.92\linewidth, height=5.2cm, xmin=0, xmax=100, ymin=0, ymax=450,
  xlabel={B, \%}, ylabel={$\lambda$, отн. ед.}, xtick={0,15,85,100},
  ytick=\empty]
  \fill[fmGray!12] (axis cs:15,0) rectangle (axis cs:85,450);
  \addplot[fmBlue, thick, domain=0:15, samples=40] {400/(1+0.45*x)};
  \addplot[fmBlue, thick] coordinates {(15,51.6) (85,65.2)};
  \addplot[fmBlue, thick, domain=85:100, samples=40] {300/(1+0.24*(100-x))};
  \node[font=\scriptsize] at (axis cs:7,420) {$\alpha$};
  \node[font=\scriptsize] at (axis cs:50,200) {$\alpha+\beta$};
  \node[font=\scriptsize] at (axis cs:93,420) {$\beta$};
\end{axis}
\end{tikzpicture}
\caption{}
\end{subfigure}
\caption{Теплопроводность сплавов: \textit{а}~— непрерывный ряд твёрдых
растворов Ag--Au (оценка по закону Видемана--Франца при 300~К с
электросопротивлением по правилу Нордгейма); \textit{б}~— система с
ограниченными твёрдыми растворами $\alpha$ и $\beta$ (схема): в растворах
теплопроводность резко падает с концентрацией, в двухфазной области меняется
линейно}
\label{fig:lambda-alloys}
\end{figure}
